\documentclass[../../../include/open-logic-section]{subfiles}

\begin{document}

\olfileid{sth}{ordinals}{opps}
\olsection{تالي او حدي ترتيبي عددونه}

په ورپسې څو څپرکو کښې به ترتيبي عددونه ډېر کاروو. نو ځينې ساده
مفهومونه معرفي کول به ګټور وي.

\begin{defn}
د هر ترتيبي عدد $\alpha$ \emph{تالي} $\ordsucc{\alpha}
=\alpha \cup \{\alpha\}$ دے. وايو چې $\alpha$ \emph{تالي} ترتيبي عدد دے،
که $\ordsucc{\beta} = \alpha$ د کوم ترتيبي عدد~$\beta$ دپاره وي.
وايو چې $\alpha$ \emph{حدي} ترتيبي عدد دے، که او يوازې که $\alpha$
نۀ تش وي او نۀ تالي ترتيبي عدد وي.
\end{defn}
\noindent
لاندې نتيجه ښيي چې دا د \emph{تالي} مناسب مفهوم دے:

\begin{prop}
د هر ترتيبي عدد $\alpha$ دپاره:
\begin{enumerate}
	\item $\alpha \in \ordsucc{\alpha}$؛
	\item $\ordsucc{\alpha}$ ترتيبي عدد دے؛
	\item داسې ترتيبي عدد $\beta$ نشته چې $\alpha \in \beta \in
	\ordsucc{\alpha}$ وي.
	\end{enumerate}
\end{prop}

\begin{proof}
ښکاره ده چې $\alpha \in \alpha \cup \{\alpha\} = \ordsucc{\alpha}$ دے.
همدارنګه، $\ordsucc{\alpha}$ د ترتيبي عددونو متعدي سټ دے، نو د
\olref[basic]{corordtransitiveord} له مخې ترتيبي عدد دے.
او $\alpha \in \beta \in \ordsucc{\alpha}$ ناشوني دي، ځکه چې بيا به يا
$\beta \in \alpha$ يا $\beta = \alpha$ وي، چې له
\olref[basic]{ordordered} سره په تناقض کښې دي.
\end{proof}
\noindent
دا د نامتناهيت هاخوا استقرا د ثبوت يوې بلې بڼې ته هم جواز ورکوي:

\begin{thm}[د نامتناهيت هاخوا ساده استقرا]\ollabel{simpletransrecursion}
فرض کړئ $\phi(x)$ داسې فارمول دے چې:
\begin{enumerate}
	\item $\phi(\emptyset)$؛ او
	\item د هر ترتيبي عدد $\alpha$ دپاره، که $\phi(\alpha)$ وي، نو
	$\phi(\ordsucc{\alpha})$ وي؛ او
	\item که $\alpha$ حدي ترتيبي عدد وي او $(\forall \beta \in
	\alpha)\phi(\beta)$ وي، نو $\phi(\alpha)$ وي.
\end{enumerate}
بيا $\forall \alpha \phi(\alpha)$ دے.
\end{thm}

\begin{proof}
نقيض عکس ثابتوو. نو فرض کړئ کوم ترتيبي عدد $\lnot\phi$ لري؛
$\gamma$ تر ټولو کوچنے داسې ترتيبي عدد وټاکئ. بيا يا
$\gamma = \emptyset$ دے، يا $\gamma = \ordsucc{\alpha}$ د کوم $\alpha$
دپاره دے چې $\phi(\alpha)$ لري؛ يا $\gamma$ حدي ترتيبي عدد دے او
$(\forall
\beta \in \gamma)\phi(\beta)$ دي.
\end{proof}
\noindent
يوه وروستۍ نښه به وروسته ګټوره ثابته شي:

\begin{defn}\ollabel{defsupstrict}
که $X$ د ترتيبي عددونو سټ وي، نو $\supstrict(X) = \bigcup_{\alpha
\in X} \ordsucc{\alpha}$ دے.
\end{defn}
\noindent
دلته «lsub» د «تر ټولو لږ په کلکه پاسني حد» نښه ده.\footnote{ځينې کتابونه
دې ته «$\text{sup}(X)$» کاروي. خو نور کتابونه «$\text{sup}(X)$» د تر ټولو
لږ \emph{غير اکيد} پاسني حد دپاره کاروي، يعنې يوازې د $\bigcup X$ دپاره.
که $X$ تر ټولو لوے غړے $\alpha$ ولري، نو دا دوه مفهومونه بېلېږي:
تر ټولو لږ \emph{اکيد} پاسنے حد $\ordsucc{\alpha}$ دے، خو تر ټولو لږ
\emph{غير اکيد} پاسنے حد يوازې $\alpha$ دے.} لاندې نتيجه دا څرګندوي:

\begin{prop}
که $X$ د ترتيبي عددونو سټ وي، نو $\supstrict(X)$ هغه تر ټولو کوچنے
ترتيبي عدد دے چې د $X$ له هر ترتيبي عدد څخه لوے وي.
\end{prop}

\begin{proof}
$Y = \Setabs{\ordsucc{\alpha}}{\alpha \in X}$ واخلو، نو
$\supstrict(X) = \bigcup Y$ دے. ځکه چې ترتيبي عددونه متعدي دي او د هر
ترتيبي عدد هر غړے ترتيبي عدد دے، $\supstrict(X)$ د ترتيبي عددونو متعدي
سټ دے، نو د \olref[basic]{corordtransitiveord} له مخې ترتيبي عدد دے.

که $\alpha \in X$ وي، نو $\ordsucc{\alpha} \in Y$ دے، نو $\ordsucc{\alpha}
\subseteq \bigcup Y = \supstrict(X)$ دے، او له دې $\alpha \in
\supstrict(X)$ دے. نو $\supstrict(X)$ د $X$ له هر ترتيبي عدد څخه په کلکه
لوے دے.

په معکوسه توګه، که $\alpha \in \supstrict(X)$ وي، نو $\alpha \in
\ordsucc{\beta} \in Y$ د کوم $\beta \in X$ دپاره دي، نو $\alpha \leq
\beta \in X$ دي. په دې توګه $\supstrict(X)$ د $X$ تر ټولو \emph{لږ}
اکيد پاسنے حد دے.
\end{proof}

\end{document}
